\section*{Limitations}

Our three benchmarks cover conversational memory with images in English. MemEye has eight scenario histories, so question-level intervals overstate its effective sample size (Appendix~\ref{app:robust}); the gold rounds of MemLens are derived from session-level annotations, which makes full-set recall an imperfect proxy there; and VisualMem asks multiple-choice questions about generated images. Our scope is memory whose answers depend on images (Appendix~\ref{app:robust}). All readers, the selector, the note writer, and the captioner are Gemini models, and selectors from other families select fewer rounds with our prompt (Appendix~\ref{app:transfer}). The judges are language models from two families rather than human annotators, and the strong selector decodes nondeterministically. The memory systems run with their default settings and our shared models rather than with settings tuned for these benchmarks. Costs use the list prices and caching terms of one provider.

\section*{Ethics Statement}

A persistent memory of a personal assistant stores raw conversations and images, and \vgs{} adds derived artifacts: image descriptions, index lines, notes, and, for the cache-native index, a copy of the index held by the model provider for the lifetime of a cache. Each artifact can expose personal information more broadly than the original exchange and must follow the retention and deletion rules of the raw history, so deleting a round must also remove its index line and invalidate the cached index. Per-user isolation, access control, and consent to provider-side storage remain obligations of a deployed system. Incomplete evidence sets can yield confident but wrong counts, and \vgs{} does not detect missing members. All experiments use public research benchmarks and involve no new data from human participants.
